\chapter{Замороженный носитель и запрет автоматических трёх поколений}

Предыдущий перебор оставил multiplicity поколений \(g\) свободным. Прежде
чем вводить новую family algebra, необходимо проверить, не выбирает ли уже
замороженный носитель
\begin{equation}
K=\mathbb{RP}^3\times S^1
\end{equation}
каноническую тройку через свои простейшие топологические или спектральные
инварианты.

\section{Точная клеточная гомология}

Стандартный клеточный комплекс \(\mathbb{RP}^3\) имеет по одной клетке в
размерностях \(0,1,2,3\) и дифференциалы
\begin{equation}
d_1=0,\qquad d_2=2,\qquad d_3=0.
\end{equation}
После произведения с клеточным комплексом \(S^1\) получаем
\begin{align}
H_0(K;\mathbb Z)&=\mathbb Z,\\
H_1(K;\mathbb Z)&=\mathbb Z\oplus\mathbb Z_2,\\
H_2(K;\mathbb Z)&=\mathbb Z_2,\\
H_3(K;\mathbb Z)&=\mathbb Z,\\
H_4(K;\mathbb Z)&=\mathbb Z.
\end{align}
Универсальная коэффициентная теорема даёт
\begin{align}
H^0(K;\mathbb Z)&=\mathbb Z,&
H^1(K;\mathbb Z)&=\mathbb Z,\\
H^2(K;\mathbb Z)&=\mathbb Z_2,&
H^3(K;\mathbb Z)&=\mathbb Z\oplus\mathbb Z_2,\\
H^4(K;\mathbb Z)&=\mathbb Z.
\end{align}
Сумма свободных Betti numbers равна 4, число torsion generators равно 2, а
Euler characteristic равна 0.

\begin{remark}
Свободные Betti numbers равны \((1,1,0,1,1)\), поэтому их арифметическая
сумма также равна 4. Значение 5 не следует ни из homology, ни из cohomology.
\end{remark}

\section{Spin-структуры и плоские характеры}

Поскольку
\begin{equation}
H^1(K;\mathbb Z_2)\cong\mathbb Z_2\oplus\mathbb Z_2,
\end{equation}
носитель имеет четыре spin-структуры. Кроме того,
\begin{equation}
\pi_1(K)=\mathbb Z_2\times\mathbb Z,
\end{equation}
поэтому пространство плоских unitary characters имеет вид
\begin{equation}
\operatorname{Hom}(\pi_1(K),U(1))\cong\{+1,-1\}\times U(1).
\end{equation}
Оно содержит две дискретные torsion-ветви и непрерывную окружность, но не
канонический трёхэлементный набор.

\section{Dirac zero-mode gate}

Для круглого \(\mathbb{RP}^3\) радиуса \(R\) scalar curvature равна
\begin{equation}
\mathcal R_K=\frac{6}{R^2}>0.
\end{equation}
Для любого flat unitary twist формула Лихнеровича даёт
\begin{equation}
D_{K,\rho}^2=\nabla_\rho^*\nabla_\rho+\frac{\mathcal R_K}{4}
\geq \frac{3}{2R^2}.
\end{equation}
Следовательно,
\begin{equation}
\ker D_{K,\rho}=0
\end{equation}
для всех четырёх spin-структур и всех плоских characters. Семейная тройка не
может быть multiplicity нулевых спинорных мод этого frozen operator.

\begin{theorem}[Frozen-\(K\) family-selector no-go]
Ни один элемент объявленного канонического меню --- свободные Betti numbers,
torsion generators, spin structures, плоские torsion characters,
фундаментальные generators или flat-twisted Dirac kernels --- не выбирает
\(g=3\). В этом меню естественные конечные multiplicities равны 0, 2 или 4.
\end{theorem}

\begin{nogocriterion}[Запрет арифметической сборки тройки]
Запрещено объявлять \(1+2=3\), разность \(4-1=3\) или произвольный выбор
трёх элементов из четырёх spin-ветвей family selector-ом без оператора,
симметрии или variational principle, канонически реализующего соответствующую
проекцию.
\end{nogocriterion}

Результат не исключает новую family algebra, non-flat bundle, defect или
другую индексную задачу. Он доказывает, что тройка не скрыта автоматически в
простейших инвариантах уже замороженного \(K\). Машинный ledger сохранён в
\path{s2t_v4_frozen_k_family_selector_gate_results.json}.